%% Beginning of file 'sample701.tex'
%%
%% Version 7.0.1. Created May 2025.
%% Version 7. Created January 2025.  
%%
%% AASTeX v7+ calls the following external packages:
%% times, hyperref, ifthen, hyphens, longtable, xcolor, 
%% bookmarks, array, rotating, ulem, and lineno 
%%
%% RevTeX is no longer used in AASTeX v7+.
%%
\documentclass[trackchanges, twocolappendix, twocolumn]{aastex701}
%%
%% This initial command takes arguments that can be used to easily modify 
%% the output of the compiled manuscript. Any combination of arguments can be 
%% invoked like this:
%%
%% \documentclass[argument1,argument2,argument3,...]{aastex701}
%%
%% Six of the arguments are typestting options. They are:
%%
%%  twocolumn   : two text columns, 10 point font, single spaced article.
%%                This is the most compact and represent the final published
%%                derived PDF copy of the accepted manuscript from the publisher
%%  default     : one text column, 10 point font, single spaced (default).
%%  manuscript  : one text column, 12 point font, double spaced article.
%%  preprint    : one text column, 12 point font, single spaced article.  
%%  preprint2   : two text columns, 12 point font, single spaced article.
%%  modern      : a stylish, single text column, 12 point font, article with
%% 		  wider left and right margins. This uses the Daniel
%% 		  Foreman-Mackey and David Hogg design.
%%
%% Note that you can submit to the AAS Journals in any of these 6 styles.
%%
%% There are other optional arguments one can invoke to allow other stylistic
%% actions. The available options are:
%%
%%   astrosymb    : Loads Astrosymb font and define \astrocommands. 
%%   tighten      : Makes baselineskip slightly smaller, only works with 
%%                  the twocolumn substyle.
%%   times        : uses times font instead of the default.
%%   linenumbers  : turn on linenumbering. Note this is mandatory for AAS
%%                  Journal submissions and revisions.
%%   trackchanges : Shows added text in bold.
%%   longauthor   : Do not use the more compressed footnote style (default) for 
%%                  the author/collaboration/affiliations. Instead print all
%%                  affiliation information after each name. Creates a much 
%%                  longer author list but may be desirable for short 
%%                  author papers.
%% twocolappendix : make 2 column appendix.
%%   anonymous    : Do not show the authors, affiliations, acknowledgments,
%%                  and author contributions for dual anonymous review.
%%  resetfootnote : Reset footnotes to 1 in the body of the manuscript.
%%                  Useful when there are a lot of authors and affiliations
%%		    in the front matter.
%%   longbib      : Print article titles in the references. This option
%% 		    is mandatory for PSJ manuscripts.
%%
%% Since v6, AASTeX has included \hyperref support. While we have built in 
%% specific %% defaults into the classfile you can manually override them 
%% with the \hypersetup command. For example,
%%
%% \hypersetup{linkcolor=red,citecolor=green,filecolor=cyan,urlcolor=magenta}
%%
%% will change the color of the internal links to red, the links to the
%% bibliography to green, the file links to cyan, and the external links to
%% magenta. Additional information on \hyperref options can be found here:
%% https://www.tug.org/applications/hyperref/manual.html#x1-40003
%%
%% The "bookmarks" has been changed to "true" in hyperref
%% to improve the accessibility of the compiled pdf file.
%%
%% If you want to create your own macros, you can do so
%% using \newcommand. Your macros should appear before
%% the \begin{document} command.
%%
\usepackage{amsmath}
\usepackage{bm}
\usepackage{CJK}
\usepackage{comment}
\usepackage{txfonts} % replaces text and math fonts with Times-style fonts
\newcommand{\vdag}{(v)^\dagger}
\newcommand\aastex{AAS\TeX}
\newcommand\latex{La\TeX}
%%%%%%%%%%%%%%%%%%%%%%%%%%%%%%%%%%%%%%%%%%%%%%%%%%%%%%%%%%%%%%%%%%%%%%%%%%%%%%%%
%%
%% The following section outlines numerous optional output that
%% can be displayed in the front matter or as running meta-data.
%%
%% Running header information. A short title on odd pages and 
%% short author list on even pages. Note that this
%% information may be modified in production.
\shorttitle{\textsc{Arepo-CD}}
\shortauthors{C. Kong et al.}
%%
%% Include dates for submitted, revised, and accepted.
\received{\today}
\revised{\today}
\accepted{\today}
%%
%% Indicate AAS Journal the manuscript was submitted to.
\submitjournal{APJ}
%% Note that this command adds "Submitted to " the argument.
%%
%% You can add a light gray and diagonal water-mark to the first page 
%% with this command:
%% \watermark{text}
%% where "text", e.g. DRAFT, is the text to appear.  If the text is 
%% long you can control the water-mark size with:
%% \setwatermarkfontsize{dimension}
%% where dimension is any recognized LaTeX dimension, e.g. pt, in, etc.
%%%%%%%%%%%%%%%%%%%%%%%%%%%%%%%%%%%%%%%%%%%%%%%%%%%%%%%%%%%%%%%%%%%%%%%%%%%%%%%%
%%
%% Use this command to indicate a subdirectory where figures are located.
%%\graphicspath{{./}{figures/}}
%% This is the end of the preamble.  Indicate the beginning of the
%% manuscript itself with \begin{document}.

\begin{document}
\begin{CJK*}{UTF8}{gbsn}

\title{Connecting intermediate-mass black holes formed in globular clusters to high redshift supermassive black holes}

%% A significant change from AASTeX v6+ is in the author blocks. Now an email
%% address is required for each author. This means that each author requires
%% at least one of the following:
%%
%% \author
%% \affiliation
%% \email
%%
%% If these three commands are not available for each author, the latex
%% compiler will issue an error and if you force the latex compiler to continue,
%% it will generate an incomplete pdf.
%%
%% Multiple \affiliation commands are allowed and authors can also include
%% an optional \altaffiliation to indicate a status, i.e. Hubble Fellow. 
%% while affiliations are indexed as footnotes, altaffiliations are noted with
%% with a non-numeric footnote that is set away from the numeric \affiliation 
%% footnotes. NOTE that if an \altaffiliation command is used it must 
%% come BEFORE the \affiliation call, right after the \author command, in 
%% order to place the footnotes in the proper location. Because non-numeric
%% symbols are used, \altaffiliation should be used sparingly.
%%
%% In v7+ the \author command takes an optional argument which provides 
%% additional metadata about the author. Authors can provide the 16 digit 
%% ORCID, the surname (family or last) name, the given (first or fore-) name, 
%% and a name suffix, e.g. "Jr.". The syntax is:
%%
%% \author[orcid=0000-0002-9072-1121,gname=Gregory,sname=Schwarz]{Greg Schwarz}
%%
%% This name metadata in not shown, it is only for parsing by the peer review
%% system so authors can be more easily identified. This name information will
%% also be sent to the publisher so they can include it in the CROSSREF 
%% metadata. Including an orcid will hyperlink the author name to the 
%% author's ORCID page. Note that  during compilation, LaTeX will do some 
%% limited checking of the format of the ID to make sure it is valid. If 
%% the "orcid-ID.png" image file is  present or in the LaTeX pathway, the 
%% ORCID icon will appear next to the authors name.
%%
%% Even though emails are now required for each author, the \email does not
%% produce output in the compiled manuscript unless the optional "show" command
%% is used. For example,
%%
%% \email[show]{greg.schwarz@aas.org}
%%
%% All "shown" emails are show in the bottom left of the first page. Due to
%% space constraints, only a few emails should be shown. 
%%
%% To identify a corresponding author, use the \correspondingauthor command.
%% The command appends "Corresponding Author: " to the argument it appears at
%% the bottom left of the first page like the output from \email. 

\author[0009-0003-9099-5805,gname=Chuizheng,sname=Kong]{Chuizheng Kong (孔垂政)}
\affiliation{Department of Astronomy, Tsinghua Univeristy, Beijing, 100084, People's Republic of China}
\email{kcz24@mails.tsinghua.edu.cn}

\author[0000-0002-1253-2763,gname=Hui,sname=Li]{Hui Li (李辉)}
\affiliation{Department of Astronomy, Tsinghua Univeristy, Beijing, 100084, People's Republic of China}
\email[show]{hliastro@tsinghua.edu.cn}

%% Use the \collaboration command to identify collaborations. This command
%% takes an optional argument that is either a number or the word "all"
%% which tells the compiler how many of the authors above the command to
%% show. For example "\collaboration[all]{(DELVE Collaboration)}" wil include
%% all the authors above this command.
%%
%% Mark off the abstract in the ``abstract'' environment. 
\begin{abstract}

We describe the current semi-analytic nuclear-star-cluster (NSC) model used to connect globular cluster (GC) formation, GC orbital evolution, and intermediate-mass black hole (IMBH) seeding. The model paints GC formation onto cosmological halo merger trees, samples individual clusters from a Schechter initial cluster mass function, assigns birth radii from a Gao+2023 S{\'e}rsic spatial model, and evolves the clusters in an analytic, time-dependent galactic background. Stellar mass is removed by stellar evolution, continuous tidal stripping, direct tidal tearing, and dynamical friction-driven sinking. IMBHs are assigned once at cluster birth using a metallicity- and surface-density-dependent runaway-collision fit. If the stellar component of an IMBH-hosting GC is removed before the IMBH reaches the centre, the black hole is evolved as a wandering remnant until it either sinks or remains in the galaxy.

\end{abstract}

%% Keywords should appear after the \end{abstract} command. 
%% The AAS Journals now uses Unified Astronomy Thesaurus (UAT) concepts:
%% https://astrothesaurus.org
%% You will be asked to selected these concepts during the submission process
%% but this old "keyword" functionality is maintained in case authors want
%% to include these concepts in their preprints.
%%
%% You can use the \uat command to link your UAT concepts back its source.
\keywords{\uat{Galaxies}{573} --- \uat{Cosmology}{343} --- \uat{High Energy astrophysics}{739} --- \uat{Interstellar medium}{847} --- \uat{Stellar astronomy}{1583} --- \uat{Solar physics}{1476}}

%% From the front matter, we move on to the body of the paper.
%% Sections are demarcated by \section and \subsection, respectively.
%% Observe the use of the LaTeX \label
%% command after the \subsection to give a symbolic KEY to the
%% subsection for cross-referencing in a \ref command.
%% You can use LaTeX's \ref and \label commands to keep track of
%% cross-references to sections, equations, tables, and figures.
%% That way, if you change the order of any elements, LaTeX will
%% automatically renumber them.

\section{Introduction} \label{sect:intro}

\section{Method} \label{sect:method}

The NSC model follows the ex-situ channel in which GCs form during hierarchical galaxy assembly, orbit in an evolving host potential, and contribute stellar mass to the nuclear region through stellar evolution, tidal stripping, direct tidal tearing, and dynamical-friction-driven sinking. The model is based on the semi-analytic GC framework of Choksi+2018 and Gao+2023, extended with the newer Chen+Gnedin2023 baryon-retention limit, the Chen+Gnedin2024 mass--metallicity relation, the Fragione+2019 local tidal-stripping prescription, and a formation-time IMBH seed module based on the FROST-CLUSTERS runaway-collision fits. Where the code contains multiple historical branches, this section describes the latest implemented physical branch used by the maintained NSC runs.

\subsection{GC formation} \label{subsect:gc_formation}

We paint GC formation onto fixed halo merger trees. Unless explicitly restricted to the main progenitor branch for a controlled test, the model walks all retained branches of each tree. A formation event is triggered when the specific halo growth rate between a progenitor and its descendant exceeds the threshold \(p_3\),
\begin{equation}
R_{\rm m} =
\frac{M_{{\rm h},2}/M_{{\rm h},1}-1}{t_2-t_1}
> p_3 ,
\label{eq:formation_trigger}
\end{equation}
where \(M_{{\rm h},1}\) and \(M_{{\rm h},2}\) are the halo masses at cosmic times \(t_1\) and \(t_2\). The model then assigns a stellar mass to the host using the Behroozi et al. stellar mass--halo mass relation. Along each branch, the first stellar mass is drawn with scatter. Later timesteps preserve memory of the previous stellar mass by adding the change in the median relation multiplied by a lognormal fluctuation,
\begin{equation}
M_{\star,2} =
M_{\star,1}+
\left[\overline{M}_{\star}(M_{{\rm h},2},z_2)
-\overline{M}_{\star}(M_{{\rm h},1},z_1)\right]
10^{\mathcal{N}(0,\xi)} ,
\end{equation}
with
\begin{equation}
\xi(z)=0.218+0.023\frac{z}{1+z}.
\end{equation}
If the stochastic update would give an unphysical negative stellar mass, the previous stellar mass is retained.

The cold gas mass is computed from the gas-to-stellar mass ratio
\begin{equation}
\eta(M_{\star},z) \equiv \frac{M_{\rm gas}}{M_{\star}} =
\begin{cases}
6.8\left(M_{\star}/10^9\,M_{\odot}\right)^{-n_M}
\left[(1+z)/3\right]^{2.7}, & z<2,\\
6.8\left(M_{\star}/10^9\,M_{\odot}\right)^{-n_M}
\left[(1+z)/3\right]^{1.4}, & 2\le z\le 3,\\
10.2\left(M_{\star}/10^9\,M_{\odot}\right)^{-n_M}, & z>3,
\end{cases}
\label{eq:gas_fraction}
\end{equation}
where \(n_M=0.33\) for \(M_{\star}\ge 10^9\,M_{\odot}\) and \(n_M=0.19\) below this mass. A 0.3 dex lognormal scatter is applied to this gas relation.

For the baryon-retention limit, we use the newer Chen+Gnedin2023 branch implemented in the code. Defining \(f_{\star}=M_{\star}/(f_{\rm b}M_{\rm h})\) and \(f_{\rm gas}=M_{\rm gas}/(f_{\rm b}M_{\rm h})\), the sum \(f_{\star}+f_{\rm gas}\) is not allowed to exceed the accreted baryon fraction normalised by \(f_{\rm b}\),
\begin{equation}
f_{\rm in}=s\left(M_{\rm ch}(z)/M_{\rm h},2\right),
\end{equation}
where
\begin{equation}
s(x,y)=\left[1+\left(2^{y/3}-1\right)x^y\right]^{-3/y},
\end{equation}
\begin{equation}
M_{\rm ch}(z)=1.69\times10^{10}\,M_{\odot}
\exp\left[-0.63z-\ln\left(1+\exp\{(z/\beta)^{\gamma}\}\right)\right],
\end{equation}
and the current branch adopts \(z_{\rm rei}=6\), \(\gamma=15\), and
\begin{equation}
\beta=z_{\rm rei}
\left[\ln\left(1.82\times10^3\exp[-0.63z_{\rm rei}]\right)-1\right]^{-1/\gamma}.
\end{equation}
If the baryon limit is exceeded, the gas mass is reduced to \(M_{\rm gas}=(f_{\rm in}-f_{\star})f_{\rm b}M_{\rm h}\).

The total mass formed in clusters during the event is
\begin{equation}
M_{\rm GC,tot}=\frac{3\times10^{-5}p_2}{f_{\rm b}}M_{\rm gas},
\label{eq:gc_total_mass}
\end{equation}
where \(p_2\) sets the overall GC formation efficiency. This is equivalent to the \(1.8\times10^{-4}p_2M_{\rm gas}\) form for the adopted cosmic baryon fraction \(f_{\rm b}=0.167\). Events with \(M_{\rm GC,tot}<10^5\,M_{\odot}\) do not form any cluster in the current NSC implementation.

\subsection{GC sampling} \label{subsect:gc_sampling}

\subsubsection{Mass sampling} \label{subsubsect:mass_sampling}

Individual GC birth masses are sampled from a Schechter initial cluster mass function,
\begin{equation}
\frac{dN}{dM}\propto M^{-2}\exp(-M/M_{\rm c}),
\label{eq:cimf}
\end{equation}
with \(M_{\rm min}=10^5\,M_{\odot}\). The cut-off mass \(M_{\rm c}\) is an input parameter through \(\log_{10}(M_{\rm c}/M_{\odot})\); the maintained IMBH--NSC runs use \(M_{\rm c}=10^7\,M_{\odot}\). For each formation event we first reserve the most massive cluster by solving the deterministic constraint that one cluster is expected above \(M_{\rm max}\),
\begin{equation}
1=\int_{M_{\rm max}}^{\infty}\frac{dN}{dM}\,dM .
\end{equation}
The rest of the cluster population is then drawn by inverse-transform sampling of the cumulative Schechter function until the event mass budget is exhausted. If the final draw would exceed the remaining mass budget, its mass is set to the remaining mass. The cluster list is shuffled before radii are assigned, so the deterministically reserved maximum-mass cluster is not always placed at the smallest S{\'e}rsic radius.

\subsubsection{Spatial sampling} \label{subsubsect:spatial_sampling}

Main-progenitor-branch clusters are treated as in-situ objects and are assigned birth galactocentric radii from the deprojected S{\'e}rsic profile used by Gao+2023,
\begin{equation}
\rho(r)=\rho_0\left(\frac{r}{R_{\rm e}}\right)^{-p}
\exp\left[-b\left(\frac{r}{R_{\rm e}}\right)^{1/N_{\rm S}}\right],
\label{eq:sersic_density}
\end{equation}
where
\begin{equation}
p=1-\frac{0.6097}{N_{\rm S}}+\frac{0.05563}{N_{\rm S}^2},
\qquad
b=2N_{\rm S}-\frac{1}{3}+\frac{0.009876}{N_{\rm S}}.
\end{equation}
The active model still uses the Gao+2023 halo-spin radius scale,
\begin{equation}
R_{\rm e}=\frac{j_{\rm sp}h}{20H(z)R_{\rm vir}},
\end{equation}
which is the code form of the Bullock-style spin proxy for the classical \(R_{\rm e}\simeq \lambda R_{\rm vir}/\sqrt{2}\) disc scaling. The S{\'e}rsic enclosed-mass fraction is inverted through the regularised incomplete gamma function,
\begin{equation}
r_{\rm GC}=R_{\rm e}
\left[\frac{P^{-1}\{N_{\rm S}(3-p),u\}}{b}\right]^{N_{\rm S}},
\end{equation}
where \(u\) is the assigned enclosed mass fraction of a cluster within the formation event. Clusters that form on non-main branches are treated as accreted satellite clusters and are initially placed at \(0.5R_{\rm vir}\) of the descendant halo rather than in the central S{\'e}rsic component.

\subsubsection{Metallicity assignment} \label{subsubsect:metallicity_assignment}

When a formation event occurs, the host interstellar-medium metallicity is assigned from the newer Chen+Gnedin2024 mass--metallicity branch implemented in the code,
\begin{equation}
[\mathrm{Fe}/\mathrm{H}]_{\rm gal}
=0.3\log_{10}\left(\frac{M_{\star}}{10^9\,M_{\odot}}\right)
-\log_{10}(1+z)-0.5 .
\label{eq:chen_mzr}
\end{equation}
The metallicity is capped at \([\mathrm{Fe}/\mathrm{H}]=0.3\). Each cluster then receives an independent internal metallicity perturbation,
\begin{equation}
[\mathrm{Fe}/\mathrm{H}]_{\rm GC}
=[\mathrm{Fe}/\mathrm{H}]_{\rm gal}+\mathcal{N}(0,\sigma_m),
\end{equation}
with \(\sigma_m=0.3\) dex. This scattered cluster metallicity is used both for stellar-evolution bookkeeping and for the IMBH seed calculation.

\subsection{GC evolution} \label{subsect:gc_evolution}

The active evolution solver follows every formed GC in an analytic, evolving background from its formation time to \(z=0\). The halo history is read from the main-progenitor block of the fixed merger tree; the halo mass is linearly interpolated with cosmic time, and the spin normalisation is updated from the corresponding tree block. The dark-matter component is modelled with an NFW profile, the stellar component with a S{\'e}rsic background, and previously stripped model GCs contribute to the enclosed mass used in later timesteps. In the current code, the S{\'e}rsic shape used in the background-density call is fixed at \(N_{\rm S}=2.2\), while the GC birth radii can be scanned over the requested \(N_{\rm S}\) values. The event scheduler divides the full interval into 100 coarse background blocks and evolves the fastest-changing active object first within each block, with the timestep capped by the smaller of the mass-loss and orbital-decay timescales and by the maximum solver step of \(0.1\,{\rm Gyr}\).

\subsubsection{Stellar evolution} \label{subsubsect:stellar_evolution}

Bound stellar mass is reduced by a cumulative stellar-evolution loss fraction,
\begin{equation}
f_{\rm se}(t)=
\max\left[
0,
-\frac{x^2}{100}+0.288x-1.42
\right],
\qquad
x=\log_{10}\left(\frac{t}{\rm Gyr}\right)+9 ,
\label{eq:stellar_evolution}
\end{equation}
where \(t\) is the cluster age. At each finite timestep, the increment \(\Delta f_{\rm se}\) removes the corresponding stellar mass from the bound cluster. The same stellar-evolution correction is also applied to stellar debris that was stripped at earlier timesteps, so the deposited stellar-mass profile tracks both unevolved and evolved stellar mass.

\subsubsection{Tidal stripping} \label{subsubsect:tidal_stripping}

For the continuous tidal-stripping branch we use the implemented Fragione+2019 mode, which is the latest available mode in the model. In code units \(M_5=M/(10^5\,M_{\odot})\), the stripping rate is
\begin{equation}
\dot{M}_{\rm tid}
=\frac{2^{2/3}0.1M_5^{1/3}}{P(r)},
\qquad
P(r)=100\frac{r/{\rm kpc}}{v_{\rm c}/{\rm km\,s^{-1}}},
\label{eq:tidal_stripping}
\end{equation}
where \(r\) is the instantaneous galactocentric radius and \(v_{\rm c}\) is the local circular velocity in the combined background and deposited-GC potential. The stripped stellar mass is added to the radial deposition bin occupied by the cluster at that timestep. If an IMBH is present, stripping is not allowed to remove the black-hole mass itself; only the stellar component can be lost.

\subsubsection{Direct tidal tearing} \label{subsubsect:direct_tidal_tearing}

In addition to continuous stripping, the model includes direct tidal tearing. At each active step the cluster half-mass density is approximated as
\begin{equation}
\rho_{\rm h}=
\begin{cases}
10^3, & M_5<1,\\
10^3M_5^2, & 1\le M_5\le 10,\\
10^5, & M_5>10,
\end{cases}
\label{eq:cluster_density}
\end{equation}
in the internal density units of the evolution solver. The cluster is directly destroyed if \(\rho_{\rm h}<\rho_{\rm tot}(r)\), where \(\rho_{\rm tot}\) is the sum of the analytic galactic background density and the mean enclosed density of previously deposited model GCs. A GC without an IMBH is then marked as torn and its remaining bound mass is deposited locally. A GC with an IMBH deposits its remaining stellar component, but the IMBH remnant survives as a wandering black hole.

\subsubsection{Dynamical friction} \label{subsubsect:dynamical_friction}

Dynamical friction drives the radial inspiral of both bound GCs and wandering IMBH remnants. The instantaneous inspiral rate is evaluated as
\begin{equation}
\dot{r}_{\rm df}=-\frac{M_5}{0.45\,r\,v_{\rm c}},
\label{eq:dynamical_friction}
\end{equation}
where \(M_5\) is the current inspiralling mass in units of \(10^5\,M_{\odot}\). The code integrates this radial decay with a fourth-order Runge--Kutta estimate that re-evaluates the enclosed deposited mass and analytic background density at substeps. Objects reaching \(r_{\rm sink}=10^{-3}\,{\rm kpc}\) are counted as sunk into the central NSC channel. The command-line switch \texttt{--DF} can disable this radial decay for controlled tests, but the physical model described here uses dynamical friction.

\subsection{IMBH} \label{subsect:imbh}

\subsubsection{Runaway collisions and IMBH formation} \label{subsubsect:runaway_imbh}

The IMBH module assigns a seed once, at GC birth. It follows the FROST-CLUSTERS runaway-collision prescription, in which dense, metal-poor clusters can build very massive stars through stellar collisions before those products collapse into IMBHs. The cluster structural scale is the three-dimensional half-mass radius
\begin{equation}
r_{\rm h}=
\frac{f_{\rm h}R_4}{1.3}
\left(\frac{M_{\rm cl}}{10^4\,M_{\odot}}\right)^{\beta},
\label{eq:imbh_radius}
\end{equation}
with \(f_{\rm h}=0.125\), \(R_4=2.365\,{\rm pc}\), and \(\beta=0.180\). Assuming a Plummer profile, the projected half-mass radius is \(R_{\rm h}=r_{\rm h}/1.305\), and the projected half-mass surface density is
\begin{equation}
\Sigma_{\rm h}=\frac{M_{\rm cl}}{2\pi R_{\rm h}^2}.
\label{eq:imbh_sigma}
\end{equation}

The metallicity input is \(Z/Z_{\odot}=10^{[\mathrm{Fe}/\mathrm{H}]_{\rm GC}}\), clipped to \(10^{-2}\le Z/Z_{\odot}\le 1\). The seed mass is then evaluated from two metallicity-dependent fits. Defining \(Q=\log_{10}[\Sigma_{\rm h}/(M_{\odot}\,{\rm pc}^{-2})]\), the calibrated-density branch is
\begin{equation}
M_{\bullet,9}=
\begin{cases}
A(Z)Q+B(Z)Q^2+C(Z),
& Q\ge Q_{\rm crit}(Z),\\
0, & Q<Q_{\rm crit}(Z),
\end{cases}
\end{equation}
where \(A\), \(B\), \(C\), and \(Q_{\rm crit}=\log_{10}[\Sigma_{\rm crit}/(M_{\odot}\,{\rm pc}^{-2})]\) are piecewise-linear functions of \(\log_{10}(Z/Z_{\odot})\) with the FROST-CLUSTERS coefficient intervals. For high surface densities, \(Q\ge 5.22\), the code switches to the high-density extrapolation
\begin{equation}
M_{\bullet,10}=D(Z)Q+E(Z).
\end{equation}
The adopted seed mass is \(M_{\bullet,10}\) above this density threshold and \(M_{\bullet,9}\) below it. Any seed with \(M_{\bullet}<100\,M_{\odot}\) is discarded. After assignment, the IMBH mass and the birth structural diagnostics \((r_{\rm h},\Sigma_{\rm h})\) are stored with the GC and are not updated by later stellar evolution or tidal stripping.

\subsubsection{Wandered IMBHs} \label{subsubsect:wandered_imbh}

An IMBH-hosting GC can lose its stellar envelope before reaching the centre. This happens either when continuous stellar and tidal mass loss reduces the bound stellar mass down to the IMBH mass, or when direct tidal tearing removes the remaining stellar component. In that case the model deposits the residual stellar mass into the local radial bin and converts the object into a wandering IMBH. A wandering IMBH no longer undergoes stellar evolution or tidal stripping, but it continues to experience dynamical friction using its own mass. If it reaches \(r_{\rm sink}\), it is counted as a sunk wandering IMBH; otherwise it remains as a surviving off-centre remnant at the final redshift.

%% Please use the acknowledgment and contribution environments. This will 
%% be anonomyized when the "anonymous" style option is used. 
\begin{acknowledgments}
We thank all the people that have made this AASTeX what it is today.  This
includes but not limited to Bob Hanisch, Chris Biemesderfer, Lee Brotzman,
Pierre Landau, Arthur Ogawa, Maxim Markevitch, Alexey Vikhlinin and Amy
Hendrickson. Also special thanks to David Hogg and Daniel Foreman-Mackey
for the new {\tt\string modern} style design. Considerable help was provided via bug
reports and hacks from numerous people including Patricio Cubillos, Alex
Drlica-Wagner, Sean Lake, Michele Bannister, Peter Williams, Jonathan
Gagne, Arthur Adams, Nicholas Wogan, Aaron Pearlman, Jeff Mangum, Mark Durre, Joel Ong, and Stephen Thorp.
\end{acknowledgments}

\begin{contribution}
%%This section gives authors the space to recognize author contributions. The text inside this environment is NOT counted towards the total word quanta. At a minimum, manuscripts are expected to include this text:

All authors contributed equally to the Terra Mater collaboration.

%% But authors are expected to provide more specific details, e.g. 
%%
%%SC was responsible for writing and submitting the manuscript.
%%WWM came up with the initial research concept and edited the manuscript.
%%OTS obtained the funding and edited the manuscript.
%%EBF provided the formal analysis and validation. He also edited the manuscript.
%%GEH Supervised the undergraduates, wrote the software and administers the project github and Zenodo repositories.
%%
%% Authors can use the Contributor Role Taxonomy (CRediT) at
%% https://credit.niso.org
%% for ideas on how write a good statement tailored to their needs.

\end{contribution}

%% To help institutions obtain information on the effectiveness of their 
%% telescopes the AAS Journals has created a group of keywords for telescope 
%% facilities.
%
%% Following the acknowledgments section, use the following syntax and the
%% \facility{} or \facilities{} macros to list the keywords of facilities used 
%% in the research for the paper.  Each keyword is check against the master 
%% list during copy editing.  Individual instruments can be provided in 
%% parentheses, after the keyword, but they are not verified.
\facilities{HST(STIS), Swift(XRT and UVOT), AAVSO, CTIO:1.3m, CTIO:1.5m, CXO}

%% Similar to \facility{}, there is the optional \software command to allow 
%% authors a place to specify which programs were used during the creation of 
%% the manuscript. Authors should list each code and include either a
%% citation or url to the code inside ()s when available.
\software{astropy \citep{2013A&A...558A..33A,2018AJ....156..123A,2022ApJ...935..167A},  
          Cloudy \citep{2013RMxAA..49..137F}, 
          Source Extractor \citep{1996A&AS..117..393B}
          }

%% Appendix material should be preceded with a single \appendix command.
%% There should be a \section command for each appendix. Mark appendix
%% subsections with the same markup you use in the main body of the paper.
%%
%% Each Appendix (indicated with \section) will be lettered A, B, C, etc.
%% The equation counter will reset when it encounters the \appendix
%% command and will number appendix equations (A1), (A2), etc. The
%% Figure and Table counter will not reset.

\appendix

\section{Determining halo parameters}

%% For this sample we use BibTeX plus aasjournalv7.bst to generate the
%% the bibliography. The sample7.bib file was populated from ADS. To
%% get the citations to show in the compiled file do the following:
%%
%% pdflatex sample7.tex
%% bibtext sample7
%% pdflatex sample7.tex
%% pdflatex sample7.tex

\bibliography{ref}{}
\bibliographystyle{aasjournalv7}

%% This command is needed to show the entire author+affiliation list when
%% the collaboration and author truncation commands are used.  It has to
%% go at the end of the manuscript.
%\allauthors

%% Include this line if you are using the \added, \replaced, \deleted
%% commands to see a summary list of all changes at the end of the article.
%\listofchanges

\end{CJK*}
\end{document}

% End of file `sample7.tex'.
